%=============================================================================!
% 8.F  SOLUTIONS TO THE EXERCISES                                             !
%=============================================================================!
\unnumbered{Chapter~\ref{ch:potentials}: Potential energy surfaces and
interatomic models}
\label{sec:pe-solutions}

\solutionto{ex:pe-ratio}With the same stiffness, $\omega_0 = \sqrt{k/m}$
is inversely proportional to $\sqrt{m}$, so the ratio is $\sqrt{6.94/
\num{5.486e-4}} = \sqrt{12\,650} = 112.5$.

\solutionto{ex:pe-quantum}$T = \hbar\omega/\kB$, with $\kB =
\SI{8.617e-5}{\electronvolt\per\kelvin}$. For the hollow $\hbar\omega =
0.6582\times2\pi/170.3 = \SI{0.02428}{\electronvolt}$, so $T =
0.02428/\num{8.617e-5} = \SI{282}{\kelvin}$, just below room
temperature; for the O--H stretch $T = 0.4534/\num{8.617e-5} =
\SI{5262}{\kelvin}$.

\solutionto{ex:pe-lj-pull}The pull is $\varphi'(r) = 4\varepsilon
(-12\sigma^{12}/r^{13} + 6\sigma^6/r^7)$, by the power rule for negative
powers (\cref{sec:os-anharmonic}). It is greatest where its own slope,
$\varphi''(r) = 4\varepsilon(156\sigma^{12}/r^{14} - 42\sigma^6/r^8)$, is
zero, that is where $156\sigma^{12}/r^{14} = 42\sigma^6/r^8$; multiplying
by $r^{14}/(42\sigma^6)$ gives $r^6 = (156/42)\sigma^6 = (26/7)\sigma^6$,
so $r = (26/7)^{1/6}\sigma = 1.2445\,\sigma$. There $(\sigma/r)^6 = 7/26$
and $(\sigma/r)^{12} = 49/676$, so
\begin{equation*}
   \varphi' = \frac{4\varepsilon}{r}\Bigl(-12\times\frac{49}{676}
   + 6\times\frac{7}{26}\Bigr) = \frac{4\varepsilon}{1.2445\,\sigma}\,
   (-0.8698 + 1.6154) = 2.396\,\frac{\varepsilon}{\sigma} ,
\end{equation*}
a pull of $2.396\,\varepsilon/\sigma$ towards the other atom.

\solutionto{ex:pe-morse}With $y = r - r_0$ and the chain rule, $\varphi'
= D(-2a\,\mathrm{e}^{-2ay} + 2a\,\mathrm{e}^{-ay}) = 2aD\,\mathrm{e}^{-ay}
(1 - \mathrm{e}^{-ay})$, zero only where $1 - \mathrm{e}^{-ay} = 0$, since $\mathrm{e}^{-ay}$ is
never zero, that is at $y = 0$, where $\varphi = D(1 - 2) = -D$. Differentiating again, $\varphi'' =
D(4a^2\mathrm{e}^{-2ay} - 2a^2\mathrm{e}^{-ay})$, which at $y = 0$ is
$2Da^2 > 0$: a minimum, with that stiffness.

\solutionto{ex:pe-exp6}$\varphi' = -(A/\rho)\,\mathrm{e}^{-r/\rho} +
6C/r^7$. At $r = r_{\mathrm m}$, $r/\rho = \alpha$, so $A\,\mathrm{e}^
{-\alpha} = 6\varepsilon/(\alpha - 6)$, and
\begin{equation*}
   \varphi'(r_{\mathrm m}) = -\frac{\alpha}{r_{\mathrm m}}\,
   \frac{6\varepsilon}{\alpha - 6} + \frac{6}{r_{\mathrm m}^7}\,
   \frac{\alpha\varepsilon r_{\mathrm m}^6}{\alpha - 6} = 0 ,
   \qquad
   \varphi(r_{\mathrm m}) = \frac{6\varepsilon}{\alpha - 6}
   - \frac{\alpha\varepsilon}{\alpha - 6} = -\varepsilon .
\end{equation*}
Differentiating again, $\varphi'' = (A/\rho^2)\,\mathrm{e}^{-r/\rho} -
42C/r^8$, which at $r_{\mathrm m}$ is
\begin{equation*}
   \frac{\alpha^2}{r_{\mathrm m}^2}\,\frac{6\varepsilon}{\alpha - 6}
   - \frac{42\alpha\varepsilon}{(\alpha - 6)\,r_{\mathrm m}^2}
   = \frac{6\alpha(\alpha - 7)\,\varepsilon}{(\alpha - 6)\,r_{\mathrm m}^2}
   > 0 \quad\text{for } \alpha > 7 :
\end{equation*}
a minimum.

\solutionto{ex:pe-argon}With $k = 72\varepsilon/r_{\mathrm m}^2$
(\cref{sec:pe-pairs}) and the reduced mass $m/2$, $\omega_0 =
\sqrt{2k/m} = \sqrt{144\,\varepsilon/(m r_{\mathrm m}^2)} =
(12/r_{\mathrm m})\sqrt{\varepsilon/m}$, and \cref{eq:pe-tau} gives
$\sqrt{\varepsilon/m} = \sigma/\tau$, so $\omega_0 = 12\sigma/(r_{\mathrm
m}\tau) = 12/(2^{1/6}\tau) = 10.69/\tau$. The period is $2\pi/10.69 =
0.5877\tau$, which for argon is $0.5877\times2151 =
\SI{1264}{\femto\second}$, and it spans $0.5877/0.005 = 117.5$, about 118
steps.

\solutionto{ex:pe-third}In $\sum_i\vb{F}_i = \sum_i\sum_{j\ne
i}\varphi'(r_{ij})\,\vb{r}_{ij}/r_{ij}$ each pair appears twice, as $(i,
j)$ and as $(j, i)$. The distances are equal and $\vb{r}_{ji} =
-\vb{r}_{ij}$, so the two terms cancel, and the sum is zero.

\solutionto{ex:pe-virial}$\varphi'(1.5\sigma) = 4\varepsilon(-12/1.5^{13}
+ 6/1.5^7)/\sigma = 4\varepsilon(-0.0617 + 0.3512)/\sigma =
1.158\,\varepsilon/\sigma$, so $\mathcal{W} = -1.5\times1.158 =
-1.737\varepsilon$. $\mathcal{W} = -r\varphi'(r)$ is positive where $\varphi' < 0$,
for $r < r_{\mathrm m}$, where the pair repels.

\solutionto{ex:pe-best-step}$h = (3\epsilon_{\mathrm m}\lvert U\rvert/
\lvert U'''\rvert)^{1/3} = (3\times\num{2.2e-16}\times10.75/10^4)^{1/3} =
\num{8.9e-7}\,\sigma$, close to the $10^{-6}\sigma$ measured.

\solutionto{ex:pe-forward}By the Taylor series, $U(x + h) = U + hU' +
\frac12 h^2U'' + \order{h^3}$, so $[U(x + h) - U(x)]/h = U' + \frac12 hU''
+ \order{h^2}$, an error of $\frac12 h\lvert U''\rvert$, to which rounding
adds $\epsilon_{\mathrm m}\lvert U\rvert/h$ as for central differences.
The slope of the total with respect to $h$, $\frac12\lvert U''\rvert -
\epsilon_{\mathrm m}\lvert U\rvert/h^2$, is zero at $h = \sqrt{2
\epsilon_{\mathrm m}\lvert U\rvert/\lvert U''\rvert}$. If $U$ changes by
about its own size over a length $L$, so that $\lvert U''\rvert \approx
\lvert U\rvert/L^2$ and the forces are of size $F \approx \lvert U\rvert/
L$, this step is $\sqrt{2\epsilon_{\mathrm m}}\,L$, of order
$\sqrt{\epsilon_{\mathrm m}} = \num{1.5e-8}$ times $L$, and there each
error is $\frac12 h\lvert U''\rvert = \sqrt{\epsilon_{\mathrm m}/2}\,F$,
about $10^{-8}F$. For central differences, $\lvert U'''\rvert \approx
\lvert U\rvert/L^3$ and $h \approx \epsilon_{\mathrm m}^{1/3}L$ in
\cref{eq:pe-fd-error} give an error of about $\epsilon_{\mathrm m}^{2/3}F
\approx \num{4e-11}\,F$.

\solutionto{ex:pe-switch}$S(0) = 1$ and $S(1) = 1 - 10 + 15 - 6 = 0$.
$S'(x) = -30x^2 + 60x^3 - 30x^4 = -30x^2(1 - 2x + x^2) = -30x^2(1 -
x)^2$, zero at both ends,
and $S''(x) = -60x + 180x^2 - 120x^3 = -60x(1 - x)(1 - 2x)$, zero at both
ends too, as expanding $(1 - x)(1 - 2x) = 1 - 3x + 2x^2$ confirms.

\solutionto{ex:pe-envelope}$u(1) = 1 - 28 + 48 - 21 = 0$. Differentiating
term by term, $u'(x) = -168x^5 + 336x^6 - 168x^7$ and $u''(x) = -840x^4 +
2016x^5 - 1176x^6$, so $u'(1) = -168 + 336 - 168 = 0$ and $u''(1) = -840
+ 2016 - 1176 = 0$.

\solutionto{ex:pe-tail}The atoms between $r$ and $r + \dd r$ number
$4\pi r^2n\,\dd r$, and each pair is shared by two atoms, so the energy
left out per atom is
\begin{align*}
   \frac12\int_{r_{\mathrm c}}^{\infty}4\pi r^2n\,\varphi_{\mathrm{LJ}}(r)
   \,\dd r
   &= 8\pi n\varepsilon\int_{r_{\mathrm c}}^{\infty}\Bigl(
   \frac{\sigma^{12}}{r^{10}} - \frac{\sigma^6}{r^4}\Bigr)\dd r\\
   &= 8\pi n\varepsilon\Bigl(\frac{\sigma^{12}}{9r_{\mathrm c}^9}
   - \frac{\sigma^6}{3r_{\mathrm c}^3}\Bigr)\\
   &= 8\pi n\varepsilon\sigma^3\Bigl[\frac19\Bigl(\frac{\sigma}
   {r_{\mathrm c}}\Bigr)^9 - \frac13\Bigl(\frac{\sigma}{r_{\mathrm c}}
   \Bigr)^3\Bigr]
   && \text{taking out $\sigma^3$}\\
   &= \frac83\pi n\varepsilon\sigma^3\Bigl[\frac13\Bigl(\frac{\sigma}
   {r_{\mathrm c}}\Bigr)^9 - \Bigl(\frac{\sigma}{r_{\mathrm c}}\Bigr)^3
   \Bigr] ,
\end{align*}
since $r^{-9}/(-9)$ has the derivative $r^{-10}$ and falls to zero far
out, so $\int_{r_{\mathrm c}}^\infty r^{-10}\,\dd r = r_{\mathrm c}^{-9}/
9$, and likewise $\int_{r_{\mathrm c}}^\infty r^{-4}\,\dd r = r_{\mathrm
c}^{-3}/3$. With $n = 0.8\sigma^{-3}$ and $\sigma/r_{\mathrm c} = 0.4$,
$\frac83\pi\times0.8\times(\frac13\times0.000262 - 0.064) =
-0.428\varepsilon$.

\solutionto{ex:pe-crossings}The motion is that of the shifted potential,
whose total energy is conserved, and the truncated total differs from it
by $\varphi(r_{\mathrm c}) = -0.0163\varepsilon$ for each pair inside the
cutoff. A rise of $0.49\varepsilon$ is therefore $0.49/0.0163 = 30$ pairs
fewer inside than at the start.

\solutionto{ex:pe-sqrt-force}By the chain rule, $\grad_k(-\sqrt{\rho_i}) =
-\frac12\rho_i^{-1/2}\grad_k\rho_i$, and the force is minus this.
$\rho_i$ depends on $\vb{r}_k$ only through its term $j = k$, so, with
$\grad_k r_{ik} = (\vb{r}_k - \vb{r}_i)/r_{ik}$ (\cref{sec:ro-system},
with $k$ in the place of $i$ and $i$ in the place of $j$),
\begin{equation*}
   \grad_k\rho_i = -\frac{2\zeta}{r_0}\,\xi^2\,\mathrm{e}^{-2\zeta(r_{ik}/r_0 - 1)}
   \,\frac{\vb{r}_k - \vb{r}_i}{r_{ik}} ,
\end{equation*}
which points from $k$ towards $i$: the term pulls $k$ towards $i$ with a
strength divided by $\sqrt{\rho_i}$, weaker when $i$ has many neighbours.

\solutionto{ex:pe-vacancy}With a pair potential and bonds of energy
$\varphi < 0$, each atom has $z$ bonds, each shared with one other atom,
so the energy holding one atom is $\frac12 z\lvert\varphi\rvert$. Removing
an atom breaks $z$ bonds, costing $z\lvert\varphi\rvert$; putting it back
gains $\frac12 z\lvert\varphi\rvert$; the vacancy costs $\frac12 z\lvert
\varphi\rvert$, the same. In the second-moment model each atom has the
bond energy $-\xi\sqrt{z}$, so the energy holding one atom is $\xi\sqrt
z$. Removing an atom removes its own term, costing $\xi\sqrt z$, and each
of its $z$ neighbours goes from $-\xi\sqrt z$ to $-\xi\sqrt{z - 1}$,
costing $z\xi(\sqrt z - \sqrt{z - 1})$; putting the atom back gains $\xi
\sqrt z$. The vacancy costs $z\xi(\sqrt z - \sqrt{z - 1})$, which divided
by $\xi\sqrt z$ is $z(\sqrt z - \sqrt{z - 1})/\sqrt z = z(1 - \sqrt{(z -
1)/z}) = z(1 - \sqrt{1 - 1/z})$: for $z = 12$, $12(1 -
\sqrt{11/12}) = 0.511$.

\solutionto{ex:pe-dihedral}$\vb{b}_1 = (-1, 0, 0)$, $\vb{b}_2 = (0, 0,
1)$, $\vb{b}_3 = (0, 1, 0)$. Then $\vb{n}_1 = \vb{b}_1\times\vb{b}_2 = (0,
1, 0)$ and $\vb{n}_2 = \vb{b}_2\times\vb{b}_3 = (-1, 0, 0)$, whose dot
product is 0, so $\psi = \pm\ang{90}$; $\vb{n}_1\times\vb{n}_2 = (0, 0,
1)$ points along $\vb{b}_2$, so $\psi = +\ang{90}$. There $\cos\psi = 0$,
$\cos2\psi = -1$ and $\cos3\psi = 0$, so the brackets of
\cref{eq:pe-opls} are $1 + \cos\psi = 1$, $1 - \cos2\psi = 2$ and $1 +
\cos3\psi = 1$, and the energy is $\frac12\times0.1 +
\frac12\times0.02\times2 + \frac12\times0.05 =
\SI{0.095}{\electronvolt}$.

\solutionto{ex:pe-bjerrum}$k_{\mathrm e}/r = \kB T$ gives $r = 14.3996/0.02585 =
\SI{557}{\angstrom}$.

\solutionto{ex:pe-shells}At $\sqrt3\,d$: the 8 sites $(\pm1, \pm1,
\pm1)$, with $a + b + c$ odd, so opposite in charge. At $2d$: the 6 sites
such as $(\pm2, 0, 0)$, with $a + b + c$ even, so the same charge. Each
opposite ion adds $+1/(\text{distance}/d)$ to $\mathcal{M}$ and each like ion
subtracts that: up to $d$, $\mathcal{M} = 6$; up to $\sqrt2\,d$, $6 - 12/\sqrt2 =
-2.485$; up to $\sqrt3\,d$, $-2.485 + 8/\sqrt3 = 2.134$; up to $2d$, $2.134
- 6/2 = -0.866$.

\solutionto{ex:pe-cube}The cube of half-side $d$ holds the 6 face ions at
$d$ (opposite, weight $\frac12$), the 12 edge ions at $\sqrt2\,d$ (like,
weight $\frac14$) and the 8 corner ions at $\sqrt3\,d$ (opposite, weight
$\frac18$): $\mathcal{M} = 6\times\frac12 - 12\times\frac14/\sqrt2 + 8\times
\frac18/\sqrt3 = 3 - 2.121 + 0.577 = 1.456$, already within 17\% of
1.7476.
